\begin{lemma}[Strongly closed witness implies elementarity]
Let $G=(V,E)$ be a finite loopless multigraph with
$\chi'(G)=k+1$ and $k\geq\Delta(G)+1$.  Suppose that an edge identity
$e_0$, a proper partial $k$-edge-coloring $\varphi$ leaving exactly $e_0$
uncolored, and a set $X\subseteq V$ have the following properties: $X$
contains both ends of $e_0$, and $X$ is elementary and strongly closed with
respect to $\varphi$.  Then $G$ is elementary.
\end{lemma}
\begin{proof}
We use missing colors, internal and boundary edge identities, and the two
colored-set properties with exactly the meanings in
\noderef{MissingColors}, \noderef{InternalEdges}, \noderef{BoundaryEdges},
\noderef{ColoredVertexSetElementary}, and
\noderef{ColoredVertexSetStronglyClosed}.  Properness and the assertion that
only $e_0$ is uncolored have the meanings in
\noderef{IsProperEdgeColoring} and \noderef{IsOneEdgeUncolored}.

Choose an end $x$ of $e_0$.  Such an end exists because an unordered pair
has two entries, and the loopless law in \noderef{LooplessMultigraph} says
that its two entries are distinct.  Every colored edge incident with $x$
has a different color by properness.  Since $e_0$ is incident with $x$ and
is the only uncolored edge, precisely $d_G(x)-1$ colors are present at $x$.
The definitions in \noderef{Degree}, \noderef{PresentColors}, and
\noderef{MissingColors} therefore give
\[
 |\overline\varphi(x)|=k-(d_G(x)-1)
 \geq k-(\Delta(G)-1)\geq2,
\]
where $d_G(x)\leq\Delta(G)$ follows from \noderef{MaximumDegree}.  Choose a
missing color $\alpha\in\overline\varphi(x)$.

For each color $c$, let $m_c$ be the number of vertices of $X$ at which
$c$ is missing, let $a_c$ be the number of $c$-colored edge identities
internal to $X$, and let $b_c$ be the number of $c$-colored edge identities
in $\partial_G(X)$.  We first establish the incidence identity
\[
                 2a_c+b_c=|X|-m_c.                 \tag{1}
\]
Indeed, properness implies that a vertex at which $c$ is present is
incident with exactly one $c$-colored edge, whereas a vertex at which $c$
is missing is incident with none.  Thus the right side counts incidences of
$c$-colored edges with vertices in $X$.  Every internal edge identity has
two distinct ends in $X$ and contributes two incidences; every boundary
edge identity has exactly one end in $X$ and contributes one.  These are all
possibilities, proving (1).  This count treats parallel edges as distinct
because it counts elements of the edge type.

The chosen color $\alpha$ is missing at $x\in X$.  Elementariness of $X$
says that it cannot be missing at any other vertex of $X$, so
$m_\alpha=1$.  Strong closure includes closure, and closure says that no
boundary edge has a color missing on $X$; hence $b_\alpha=0$.  Equation
(1) now yields
\[
                      2a_\alpha=|X|-1.
\]
Consequently $|X|$ is odd.  The set $X$ contains the two distinct ends of
$e_0$, so its odd cardinality is at least three.  Thus $X$ is one of the
eligible sets in \noderef{Density}.

Let $M=\bigcup_{v\in X}\overline\varphi(v)$ and put $q=|M|$.  Since $X$
is elementary, its missing-color sets are pairwise disjoint, and therefore
\[
              q=\sum_{v\in X}|\overline\varphi(v)|. \tag{2}
\]
At a vertex $v\in X$, properness and the fact that only $e_0$ is uncolored
show that the number of present colors equals the number of colored edge
identities incident with $v$.  Hence the summand in (2) is $k$ minus that
colored degree.  Summing over $X$ counts each colored internal edge twice
and each colored boundary edge once.  The edge $e_0$ is internal because
both its ends lie in $X$, and it is the only uncolored edge; in particular
every boundary edge is colored.  It follows that
\[
 q=k|X|-2\bigl|E_{G-e_0}(X)\bigr|-|\partial_G(X)|. \tag{3}
\]
Here $E_{G-e_0}(X)$ denotes the internal edge identities other than $e_0$,
all of which are colored.

Strong closure has two effects on the boundary.  Closure excludes all $q$
colors in $M$ from boundary edges, and the remaining clause says that
distinct boundary edge identities have distinct colors.  Thus the boundary
is injected into the other $k-q$ colors, giving
\[
                       |\partial_G(X)|\leq k-q.       \tag{4}
\]
Solving (3) for the internal-edge term and applying (4), we obtain
\[
 2\bigl|E_{G-e_0}(X)\bigr|
   =k|X|-|\partial_G(X)|-q
   \geq k(|X|-1).
\]
Since $e_0$ is an additional internal edge identity,
\[
 2|E_G(X)|=2\bigl|E_{G-e_0}(X)\bigr|+2
          >k(|X|-1).                                  \tag{5}
\]
By \noderef{DensityValue} and the maximal definition in
\noderef{Density}, (5) implies $\Gamma(G)>k$.  Hence the natural ceiling in
\noderef{DensityCeil} satisfies
$\lceil\Gamma(G)\rceil\geq k+1=\chi'(G)$.  Conversely,
\noderef{DensityLowerBound} gives
$\lceil\Gamma(G)\rceil\leq\chi'(G)$.  Equality follows, and the definition
in \noderef{Elementary} says exactly that $G$ is elementary.
\end{proof}
